\chapter{Maximum principles and entire graphical flows}\label{sec:entire-graphical-flow}
The periodic estimates avoid spatial infinity. We now give an explicit maximum principle sufficient for entire graphs with controlled growth, then identify the additional existence input from \cite{EH89}, Theorem 4.6. Smoothness on a noncompact time slice alone is not a global bound.

\begin{lemma}[Normal material heat equals the graph nondivergence operator]\label{lem:graphical-heat-coordinate-cancellation}
For an entire graphical mean curvature flow, let
\[
 a^{ij}=\delta_{ij}-\frac{u_i u_j}{1+|Du|^2}.
\]
For any geometric scalar $f$ written in graph coordinates,
\[
 (\partial_t^\perp-\Delta_M)f
       =\partial_t^{\rm graph}f-a^{ij}\partial_i\partial_j f.
\]
The cancellation uses the graph equation and does not assert equality of the two time derivatives separately.
\end{lemma}
\begin{proof}
The graph metric has Christoffel symbols $\Gamma^k_{ij}=u_{ij}u_k/v^2$, as follows by projecting $F_{ij}=u_{ij}e$ onto the tangent frame. Thus
$\Delta_M f=a^{ij}f_{ij}-a^{ij}u_{ij}u_kf_k/v^2$.
The graph equation gives $a^{ij}u_{ij}=u_t$. Meanwhile the tangential velocity in Lemma \ref{lem:graphical-material-drift} has coordinate components $T^k=u_tu_k/v^2$, since the projection of $e$ has components $u_k/v^2$. Subtract the Laplacian from $\partial_t^\perp f=\partial_t^{\rm graph}f-T^kf_k$. The two first-order terms cancel. All statements are local, so the same calculation holds on a torus as well.
\end{proof}

\begin{lemma}[A maximum principle with an explicit polynomial growth condition]\label{lem:entire-polynomial-maximum-principle}
On $\mathbb R^n\times[s,T]$, let $a(x,t)$ be symmetric with $0\leq a\leq A_0 I$, and let $|b(x,t)|\leq B_0$. Suppose $f$ is continuous up to $t=s$, twice differentiable in space and once in time for $t>s$, and satisfies
\[
 f_t-a^{ij}f_{ij}-b^if_i\leq0,\qquad
 f(x,s)\leq M,
 \qquad f(x,t)\leq C_0(1+|x|^2)^m
\]
for some finite $A_0,B_0,C_0,m\geq0$, uniformly on the slab. Then $f\leq M$ throughout. Strict uniform ellipticity is not needed for this upper bound.
\end{lemma}
\begin{proof}
Choose $q>\max\{m,1\}$ and put $\psi=(1+|x|^2)^q$. Direct differentiation gives
$|D\psi|\leq q\psi$ and $|D^2\psi|\leq4q^2\psi$ in operator norm. For $K=1+4nA_0q^2+B_0q$, the function $B(x,t)=e^{K(t-s)}\psi(x)$ satisfies
$B_t-a^{ij}B_{ij}-b^iB_i\geq B>0$.
For each $\eta>0$, the function $f-M-\eta B$ is negative initially and tends to minus infinity at spatial infinity uniformly in time, because $q>m$. If it were positive somewhere, it would have a positive maximum on a compact space-time cylinder, at a time later than $s$. At that point its spatial gradient is zero, its Hessian is nonpositive and its time derivative is nonnegative (one-sided at the final time). Positive semidefiniteness of $a$ would make its parabolic operator nonnegative. Its strict subsolution inequality makes the same operator negative, a contradiction. Thus $f\leq M+\eta B$; let $\eta\downarrow0$ at each fixed point. This supplies an explicit alternative to the weighted noncompact maximum principle in \cite{EH89}, Corollary 1.1, under the growth hypotheses stated here.
\end{proof}

\begin{proposition}[Slope preservation for an entire graph with controlled growth]\label{prop:entire-graph-slope-preservation}
Let $u$ be a smooth entire graphical flow on $[s,T]$. Suppose $v(x,t)=\sqrt{1+|Du|^2}$ has at most polynomial growth in $|x|$, uniformly on this slab, and $V_s=\sup_xv(x,s)<\infty$. Then $v(x,t)\leq V_s$ throughout. No a priori curvature bound is needed for this conclusion.
\end{proposition}
\begin{proof}
Proposition \ref{prop:graphical-angle-evolution} and Lemma \ref{lem:graphical-heat-coordinate-cancellation} give
$v_t-a^{ij}v_{ij}=-|A|^2v-2v^{-1}|\nabla v|^2\leq0$.
The coefficient matrix always satisfies $0<a\leq I$, independently of whether a global slope bound has yet been proved. Apply Lemma \ref{lem:entire-polynomial-maximum-principle} with zero drift and the assumed polynomial upper growth. This proves the bound without presuming a maximum is attained on the entire graph. It is a version of \cite{EH89}, Corollary 3.2 with an explicit sufficient growth class.
\end{proof}

\begin{theorem}[Entire-graph curvature smoothing with slab bounds]\label{thm:entire-graph-curvature-smoothing}
Let $u$ be a smooth entire graphical flow on $[s,T]$, and assume $v$ and $|A|$ are bounded on this entire space-time slab. With $V_s=\sup_xv(x,s)$, for $s<t\leq T$ one has
\[
 \sup_x v(x,t)^2|A(x,t)|^2\leq\sup_xv(x,s)^2|A(x,s)|^2,
\]
\[
 \bigl(1+2(t-s)|A(x,t)|^2\bigr)v(x,t)^2\leq V_s^2,
 \qquad |A(x,t)|^2\leq\frac{V_s^2-1}{2(t-s)}.
\]
This states the bounds of \cite{EH89}, Corollary 4.2 and the base case of Proposition 4.4, in a class where the maximum principle above applies directly.
\end{theorem}
\begin{proof}
Use the exact inequalities for $P=v^2|A|^2$ and $Z=v^2+2(t-s)P$ from Section \ref{sec:graphical-weighted-curvature}. Lemma \ref{lem:graphical-heat-coordinate-cancellation} writes their operator as
$\partial_t-a^{ij}\partial_{ij}+2v^{-1}(\nabla v)^k\partial_k$.
The Euclidean coordinate norm of a tangent vector is no greater than its graph-metric norm, since $g\geq I$. Therefore Lemma \ref{lem:graphical-slope-gradient} bounds the coordinate drift by $2v|A|$, uniformly on the slab. Both $P$ and $Z$ are bounded there. Apply Lemma \ref{lem:entire-polynomial-maximum-principle} with exponent zero, choosing the drift sign to match this operator. Their starting suprema are the two right-hand sides displayed above. The final curvature bound follows from $v\geq1$, exactly as in the periodic case. The assumed slab bound is used to justify the maximum principle; it is not silently supplied by smoothness on a noncompact domain.
\end{proof}

\begin{corollary}[A bounded-slope ancient entire graph is a plane under slab curvature control]\label{cor:ancient-entire-graph-rigidity}
Let $u:\mathbb R^n\times(-\infty,T)\to\mathbb R$ be a smooth graphical flow with one uniform bound for $v$ over all space and all times. Assume additionally that $|A|$ is bounded on each closed finite time slab. Then $u(x,t)=p\cdot x+c$ for fixed $p,c$. In particular no nonstationary vertical translator satisfying these hypotheses exists.
\end{corollary}
\begin{proof}
Fix $t<T$. Theorem \ref{thm:entire-graph-curvature-smoothing} applies from every earlier $s$, and its curvature numerator is bounded independently of $s$. Let $s\to-\infty$ to conclude $A(\cdot,t)=0$. Lemma \ref{lem:graph-curvature-hessian-comparison} gives $D^2u=0$, so $u(x,t)=p(t)\cdot x+c(t)$. The graph equation now has zero right-hand side for every $x$, forcing $p'(t)=0$ and $c'(t)=0$. A vertical translator would have $u=u_0+at$; hence its speed $a$ is zero. This is a parabolic rigidity deduction with the stated global bounds, not a claim about arbitrary local or strip graphs.
\end{proof}

\begin{proposition}[Uniqueness for entire flows with bounded slope and curvature]\label{prop:entire-graph-bounded-geometry-uniqueness}
Let $u_1,u_2$ be smooth entire graphical flows on $[0,T]$ with the same initial function. Suppose their slopes and curvatures are bounded on the whole slab. Then $u_1=u_2$ on this slab. The same statement holds on every finite slab of two all-time flows in this class.
\end{proposition}
\begin{proof}
Let $w=u_2-u_1$. Subtracting the two graphical equations, as in Proposition \ref{prop:graphical-mcf-comparison}, gives
\[
 w_t=a^{ij}(Du_2)w_{ij}+b^kw_k,\qquad
 b^k=\int_0^1\partial_{p_k}a^{ij}(Du_1+\theta Dw)\,d\theta\,(u_1)_{ij}.
\]
The slopes bound every interpolating gradient, and Lemma \ref{lem:graph-curvature-hessian-comparison} bounds $D^2u_1$ by its curvature and slope bounds. Hence $b$ is uniformly bounded. Also $|w(x,t)|\leq C(1+|x|)$ on the slab: each gradient is bounded, and the heights at $x=0$ are bounded by continuity on the compact time interval. Apply Lemma \ref{lem:entire-polynomial-maximum-principle} to $w$ and $-w$, both with zero initial value. This gives both inequalities and therefore equality. No global bound for the heights themselves is imposed.
\end{proof}

\begin{theorem}[Ecker--Huisken existence for globally Lipschitz entire graphs]\label{thm:ecker-huisken-global-graph-existence}
Let $u_0:\mathbb R^n\to\mathbb R$ be globally Lipschitz, $n\geq1$, with $\sqrt{1+|Du_0|^2}\leq V_0$ almost everywhere. There is a graphical mean curvature flow for all $t>0$, smooth at positive time, taking initial value $u_0$ by locally uniform convergence as $t\downarrow0$. It satisfies
\[
 v\leq V_0,\qquad t^{m+1}|\nabla^m A|^2\leq C(m,n,V_0)
 \quad(m=0,1,2,\ldots).
\]
If the initial graph is smooth with bounded covariant curvature derivatives of every order, the construction is smooth up to $t=0$ and has bounded curvature and its derivatives on each finite time slab. This is \cite{EH89}, Theorem 4.6, printed p.~464, together with Corollary 3.2 and Propositions 4.3--4.4. The parabolic existence, higher-derivative estimates and approximation arguments remain external. No existence theorem for arbitrary merely locally Lipschitz data without a global slope bound is asserted from this source.
\end{theorem}

\begin{corollary}[Smooth periodic initial graphs exist globally and converge]\label{cor:periodic-graph-global-existence-convergence}
Every smooth function on a flat torus is the initial value of a unique all-time smooth periodic graphical mean curvature flow. This solution has the convergence, area decay and total dissipation proved in Section \ref{sec:periodic-graph-convergence}. Its existence uses the external theorem above; those estimates have local proofs in the project.
\end{corollary}
\begin{proof}
Lift the initial function periodically to Euclidean space. Compactness of the torus bounds its slope and all derivatives, hence all covariant curvature derivatives of the graph. Theorem \ref{thm:ecker-huisken-global-graph-existence} gives an entire flow smooth up to time zero with bounded slope and curvature on every finite slab. Translating this solution by a lattice vector in the base gives another solution with the same initial value and bounds. Proposition \ref{prop:entire-graph-bounded-geometry-uniqueness} makes the two solutions equal. Thus the entire flow remains periodic and descends to the torus. The same uniqueness argument, or its compact maximum-principle version, gives uniqueness of the periodic solution. It is all-time smooth, so every hypothesis of the convergence and dissipation results is now satisfied. This use of a cited existence theorem is distinct from proving that theorem from the curvature estimate alone.
\end{proof}
